% 163 Exam 2C — LaTeX replica of 163Exm2C.pdf
% Compile with: xelatex main.tex  (two passes; fonts embedded from ./fonts/)
% Page geometry replicates the source PDF exactly (US Letter, same as 2A —
% NOT A4 like 1B/1C/2B — 612x792bp, coordinates in bp from top-left).
\documentclass{article}
\usepackage[paperwidth=612bp,paperheight=792bp,margin=0bp]{geometry}
\usepackage{amsmath}
\usepackage{fontspec}
\usepackage{unicode-math}
\usepackage{tikz}
\usetikzlibrary{calc}

% ---- fonts (embedded subsets extracted from the original PDF) ----
\setmainfont[Path=./fonts/, BoldFont=BookAntiqua-Bold_11.ttf]{BookAntiqua_5.ttf}
\newfontfamily\CalibriFont[Path=./fonts/]{Calibri_9.ttf}
\newfontfamily\CambriaTextFont[Path=./fonts/]{CambriaM.ttf}
% Prime (U+2032, Boolean-complement notation x'/y'/z'/w') is absent from
% EVERY Cambria Math subset extracted anywhere in this project so far (1B/
% 1C/2A/2B all lack it too -- none of those exams used BA complement
% notation, so it never got painted into any of their PDFs either, meaning
% there is no same-glyph-numbering Cambria donor available). Fallback: STIX
% Two Math's own prime glyph (open, serif math font, closest visual match
% available locally) -- deliberate, documented substitution, not silent.
\newfontfamily\PrimeFont[Path=./fonts/]{STIXTwoMath-Regular.otf}
% SymbolMT_15.ttf / ArialMT_20.ttf / TimesNewRomanPSMT_34.ttf kept in fonts/
% for provenance; loaded only if a page turns out to need them.
\newfontfamily\ArialFont[Path=./fonts/]{ArialMT_20.ttf}
\newfontfamily\TimesFont[Path=./fonts/]{TimesNewRomanPSMT_34.ttf}
\newfontfamily\HelveticaFont{Helvetica}
\setmathfont[Path=./fonts/]{CambriaM_nokern.ttf}

\pagestyle{empty}
\setlength\parindent{0pt}
\frenchspacing
\delimiterfactor=1000
\delimitershortfall=0bp
% bullet drawn manually (see 1C preamble note) — reused verbatim; re-measure
% against 2A's own subset if a page turns out to need U+2022.
\newcommand\SymBullet{\tikz[baseline=0pt]{\fill (0.224em,0.29em) circle[radius=0.174em];}}

% ---- math helpers matching the source document's Word-equation quirks ----
\newcommand\rng{.\,.}                    % the "0..9" integer-range dots
\newcommand\tcolon{\text{\CambriaTextFont :}} % ASCII colon (unicode-math maps ':' to U+2236)
\newcommand\cd{\tcolon\ {=}}             % ": =" as typed in the source
\newcommand\iw[1]{\mathit{#1}}           % italic words inside math
% script sizes as used by the source (base -> script)
\DeclareMathSizes{14.04}{14.04}{9.96}{7}
\DeclareMathSizes{12.96}{12.96}{9}{6.5}
\DeclareMathSizes{12}{12}{9}{6.5}
\DeclareMathSizes{11.04}{11.04}{8.04}{6}
\DeclareMathSizes{9.96}{9.96}{7.2}{5.4}
% invisible strut: \strutbp{height}{depth} forces delimiter sizes
\newcommand\strutbp[2]{\vrule width 0bp height #1bp depth #2bp}
% \CT{x}{y}{size}{text}: Cambria run set in TEXT mode (GPOS kerning like Word)
\newcommand\CT[4]{\TL{#1}{#2}{{\CambriaTextFont\fsz{#3}#4}}}
% \PRIME{x}{y}{size}: standalone Boolean-complement prime via the STIX
% fallback (see \PrimeFont note above) -- same placement contract as \CT.
\newcommand\PRIME[3]{\TL{#1}{#2}{{\PrimeFont\fsz{#3}\char"2032}}}
% \TDEL{x}{y-baseline}{size}{strut-h}{strut-d}{delim}: forced tall delimiter
\newcommand\TDEL[6]{\MT{#1}{#2}{#3}{\left#6\strutbp{#4}{#5}\right.}}
% \GD{x}{y-baseline}{size}{glyphid}: exact Cambria glyph (variant/assembly
% piece) at the original PDF's pen position (see tools/fixglyph.py)
\newcommand\GD[4]{\TL{#1}{#2}{{\CambriaTextFont\fsz{#3}\XeTeXglyph #4\relax}}}

% ---- size-switch helpers: \BA{13.02} -> Book Antiqua at 13.02bp ----
\newcommand\fsz[1]{\fontsize{#1bp}{#1bp}\selectfont}
\newcommand\BA[1]{\normalfont\fsz{#1}}
\newcommand\BAB[1]{\normalfont\bfseries\fsz{#1}}
\newcommand\CAL[1]{\CalibriFont\fsz{#1}}
\newcommand\CMT[1]{\CambriaTextFont\fsz{#1}}
\newcommand\SYM[1]{\normalfont\fsz{#1}}% size-only: SymbolMT runs are drawn bullets
\newcommand\HEL[1]{\HelveticaFont\fsz{#1}}
\newcommand\ARL[1]{\ArialFont\fsz{#1}}
\newcommand\TMS[1]{\TimesFont\fsz{#1}}

% ---- placement macros (coordinates in bp from page top-left) ----
% \TL{x}{y-baseline}{content}          left-anchored line
\newcommand\TL[3]{\node[anchor=base west,inner sep=0pt,outer sep=0pt]
  at ($(current page.north west)+(#1bp,-#2bp)$) {#3};}
% \JL{x}{y}{width}{content}            justified (stretched) line
\newcommand\JL[4]{\node[anchor=base west,inner sep=0pt,outer sep=0pt]
  at ($(current page.north west)+(#1bp,-#2bp)$) {\makebox[#3bp][s]{#4}};}
% \RECT{x0}{y0}{x1}{y1}                filled rule rectangle
\newcommand\RECT[4]{\fill ($(current page.north west)+(#1bp,-#2bp)$)
  rectangle ($(current page.north west)+(#3bp,-#4bp)$);}
% \MM{x}{y}{size}{math}                inline math anchored at baseline
\newcommand\MM[4]{\TL{#1}{#2}{\fsz{#3}$\displaystyle #4$}}
\newcommand\MT[4]{\TL{#1}{#2}{\fsz{#3}$#4$}} % textstyle math

% header/footer chrome for every page (measured directly from 163Exm2C.pdf
% via fitz span extraction). No "163 Exam 2C" repeating chrome exists on
% this exam (that string only appears once, on page 600's top sheet, which
% is real page content, not chrome). Footer is BookAntiqua 11.04, not bold:
% "You are free to talk to anyone" (left) + right-aligned folio number.
\newcommand\ChromePage[1]{%
  \TL{72.02}{45.6}{\BA{9.96}Due by2359 hours on 20260817}%
  \TL{72.02}{725.86}{\BA{11.04}You are free to talk to anyone}%
  \TL{520.9}{725.86}{\BA{11.04}#1}}

\newenvironment{pdfpage}{\mbox{}\begin{tikzpicture}[remember picture,overlay]}%
{\end{tikzpicture}\clearpage}


% ======== answer layer (not yet used — placeholders reused from 1C) ========
\newcommand\FOLIO[1]{\node[anchor=south east, draw, line width=0.5bp,
  inner xsep=3bp, inner ysep=2bp]
  at ($(current page.south east)+(-36bp,36bp)$) {\CMT{9.96}#1};}
\newcommand\PLBL[3]{\node[anchor=north west, draw, line width=0.5bp,
  inner xsep=2.5bp, inner ysep=2bp]
  at ($(current page.north west)+(#1bp,-#2bp)$) {\BA{9.5}#3};}
\newcommand\CIRC[4]{\draw[line width=0.7bp]
  ($(current page.north west)+(#1bp,-#2bp)$) ellipse [x radius=#3bp, y radius=#4bp];}
\newcommand\DLINE[5]{\draw[line width=#5bp]
  ($(current page.north west)+(#1bp,-#2bp)$) --
  ($(current page.north west)+(#3bp,-#4bp)$);}
\newcommand\worksheet[1]{\begin{pdfpage}\input{sheets/ws#1.tex}\FOLIO{#1}\end{pdfpage}}
